% Part: first-order-logic
% Chapter: axiomatic-deduction
% Section: proof-theoretic-notions

\documentclass[../../../include/open-logic-section]{subfiles}

\begin{document}

\iftag{FOL}
      {\olfileid{fol}{axd}{ptn}}
      {\olfileid{pl}{axd}{ptn}}

\olsection{నిరూపణ-సిద్ధాంత భావనలు}

\begin{explain}
మనం అనేక ముఖ్యమైన అర్థవిజ్ఞాన భావనలను
(\iftag{FOL}{సార్వత్రిక సత్యత్వం}{సర్వసత్యత}, అనుగమనం, సంతృప్తత)
నిర్వచించినట్లే, ఇప్పుడు వాటికి సరిపోలే \emph{నిరూపణ-సిద్ధాంత
భావనలను} నిర్వచిస్తాం. వీటిని !!{structure}sలో !!{sentence}s
సంతృప్తి చెందడాన్ని ఆశ్రయించి నిర్వచించం; కొన్ని సూత్రాల
!!{derivability} లేదా !!{nonderivability}ను ఆశ్రయించి నిర్వచిస్తాం.
ఈ భావనలు ఏకీభవిస్తాయని కనుగొనడం ఒక ముఖ్యమైన ఆవిష్కరణ. అవి
ఏకీభవిస్తాయనే విషయమే \emph{నిర్దుష్టత}, \emph{సంపూర్ణత}
సిద్ధాంతాల సారం.
\end{explain}

\begin{defn}[!!^{derivability}]
~$!A$తో ముగిసే ఒక !!a{derivation}~$\Gamma$ నుంచి ఉంటే,
!!^a{formula}~$!A$ను $\Gamma$ నుంచి \emph{!!{derivable}} అంటాం;
దీన్ని $\Gamma \Proves !A$ అని రాస్తాం.
\end{defn}

\begin{defn}[సిద్ధాంతాలు]
ఖాళీ సమితి నుంచి $!A$కు ఒక !!a{derivation} ఉంటే,
!!^a{formula}~$!A$ ఒక \emph{సిద్ధాంతం}. $!A$ సిద్ధాంతమైతే
$\Proves !A$ అని, కాకపోతే $\Proves/ !A$ అని రాస్తాం.
\end{defn}

\begin{defn}[అవైరుధ్యం]
!!{formula}s సమితి~$\Gamma$, $\Gamma\Proves/ \lfalse$ అయితే,
అప్పుడు మాత్రమే \emph{అవిరుద్ధం}; లేకపోతే అది \emph{విరుద్ధం}.
\end{defn}

\begin{prop}[స్వావర్తనత్వం]
\ollabel{prop:reflexivity}
$!A \in \Gamma$ అయితే, $\Gamma \Proves !A$.
\end{prop}

\begin{proof}
  !!{formula}~$!A$ ఒక్కటే,~$\Gamma$ నుంచి~$!A$కు ఒక
  !!a{derivation}.
\end{proof}

\begin{prop}[ఏకదిశత]
\ollabel{prop:monotonicity}
$\Gamma \subseteq \Delta$, $\Gamma \Proves !A$ అయితే,
$\Delta \Proves !A$.
\end{prop}

\begin{proof}
$\Gamma$ నుంచి $!A$కు ఉన్న ఏ !!{derivation} అయినా~$\Delta$ నుంచి
$!A$కు కూడా ఒక !!a{derivation}.
\end{proof}

\begin{prop}[సంక్రామకత్వం]
\ollabel{prop:transitivity}
$\Gamma \Proves !A$, $\{!A\} \cup \Delta \Proves !B$ అయితే,
$\Gamma \cup \Delta \Proves !B$.
\end{prop}

\begin{proof}
  $\{!A\} \cup \Delta \Proves !B$ అని అనుకుందాం. అప్పుడు
  $\{!A\} \cup \Delta$ నుంచి~$!B_1$, \dots, $!B_l = !B$ అనే ఒక
  !!a{derivation} ఉంది. ఆ !!{derivation}లోని కొన్ని దశలు, ముందరి
  పంక్తి~$!B_i = !A$ను సూచించే నియమం వల్ల సరైనవిగా ఉంటాయి.
  పరికల్పన ప్రకారం,~$\Gamma$ నుంచి~$!A$కు ఒక !!a{derivation} ఉంది;
  అంటే, ప్రతి $!A_i$ ఒక స్వీకృతమైనా,~$\Gamma$ యొక్క ఒక
  !!a{element} అయినా, నిగమన నియమం వల్ల సరైనదైనా అయ్యే
  !!a{derivation}~$!A_1$, \dots, $!A_k = !A$ ఉంది. ఇప్పుడు ఈ క్రమాన్ని
  పరిశీలించండి:
  \[
  !A_1, \dots, !A_k = !A, !B_1, \dots, !B_l = !B.
  \]
  ఇది $\Gamma \cup \Delta$ నుంచి~$!B$కు సరైన !!{derivation}; ఎందుకంటే
  ప్రతి $!B_i = !A$ ఇప్పుడు~$!A_k = !A$ను సమర్థించిన అదే నియమం వల్ల
  సమర్థించబడుతుంది.
  \sourcecorrection{OLTEAXD-001}{ఆంగ్ల మూలం ఈ చివరి వాక్యంలో
  $B_i = !A$ అని రాసి, సూత్రపు మెటాచరరాశి ముందు అవసరమైన ఆశ్చర్యార్థక
  టోకెన్‌ను వదిలింది. ముందరి నిగమన క్రమానికి సరిపోయేలా
  $!B_i = !A$గా సరిచేశాం.}
\end{proof}

దీనర్థం, ముఖ్యంగా, $\Gamma \Proves !A$, $!A \Proves !B$ అయితే
$\Gamma \Proves !B$. అలాగే, $!A_1, \dots, !A_n \Proves !B$ అయ్యి,
ప్రతి~$i$కు $\Gamma \Proves !A_i$ అయితే, $\Gamma \Proves !B$ అని
కూడా అనుసరిస్తుంది.

\begin{prop}
\ollabel{prop:incons}
ప్రతి~$!A$కు $\Gamma \Proves !A$ అయితే, అప్పుడు మాత్రమే $\Gamma$
విరుద్ధం.
\end{prop}

\begin{proof}
సాధన.
\end{proof}

\tagprob{FOL}
\begin{prob}
\olref[fol][axd][ptn]{prop:incons}ను నిరూపించండి.
\end{prob}
\tagendprob

\tagprob{notFOL}
\begin{prob}
\olref[pl][axd][ptn]{prop:incons}ను నిరూపించండి.
\end{prob}
\tagendprob

\begin{prop}[సంహతత్వం]
\ollabel{prop:proves-compact}
  \begin{enumerate}
  \item $\Gamma \Proves !A$ అయితే, $\Gamma_0 \Proves !A$ అయ్యే
    పరిమిత ఉపసమితి $\Gamma_0 \subseteq \Gamma$ ఒకటి ఉంటుంది.
  \item~$\Gamma$ యొక్క ప్రతి పరిమిత ఉపసమితి అవిరుద్ధమైతే,
    $\Gamma$ అవిరుద్ధం.
  \end{enumerate}
\end{prop}

\begin{proof}
  \begin{enumerate}
    \item $\Gamma \Proves !A$ అయితే, $!A \ident !A_n$ అయ్యేలా,
      అలాగే ప్రతి $!A_i$ ఒక తార్కిక స్వీకృతమైనా,~$\Gamma$ యొక్క ఒక
      !!a{element} అయినా, ముందరి !!{formula}s నుంచి మోడస్ పోనెన్స్
      ద్వారా వచ్చేదైనా అయ్యేలా $!A_1$, \dots,~$!A_n$ అనే పరిమిత
      !!{formula}s క్రమం ఉంటుంది. వీటిలో~$\Gamma$లో ఉన్న $!A_i$ల
      సమితినే $\Gamma_0$గా తీసుకోండి. అప్పుడు ఆ !!{derivation},
      ~ $\Gamma_0$ నుంచి కూడా ఒక !!a{derivation}; అందువల్ల
      $\Gamma_0 \Proves !A$.
    \item $!A \ident \lfalse$ అనే ప్రత్యేక సందర్భంలో ఇది~(1) యొక్క
      విపర్యయ ప్రతిజ్ఞ.
\end{enumerate}
\end{proof}

\end{document}
